\providecommand{\dascrivere}[1]{{\color{gray}\small[da scrivere: #1]}}

\noindent
Dopo quasi due anni, eccoci di nuovo al momento più complicato: trovare le parole per ringraziare chi mi ha accompagnato fin qui. Spero di non dilungarmi quanto alla cena della triennale e, soprattutto, di non trasformare anche questa cena in «un funerale», come disse mio padre. Spero anche di non ripetermi troppo, ma se alcune cose le avete già sentite, è perché in questi due anni mi avete dato ancora tanti motivi per dirvele.\par

\vspace{5mm}
\dascrivere{introduzione alla famiglia, ai genitori}\par

\vspace{5mm}
\textit{A mio padre}\par
\vspace{2mm}
\dascrivere{richiamo alla triennale (il finto duro, gli Acquario, Joe Temerario) e cosa è successo in questi due anni}\par

\vspace{5mm}
\textit{A mia madre}\par
\vspace{2mm}
\dascrivere{richiamo alla triennale (lo scricciolo che stava per laurearsi, c'eri nel mio cuore) e cosa è cambiato}\par

\vspace{5mm}
\textit{A mia sorella Giulia}\par
\vspace{2mm}
\dascrivere{dalle litigate da cane e gatto al parlarsi di tutto: come è andata dopo}\par

\vspace{5mm}
\textit{A mia sorella Flavia}\par
\vspace{2mm}
\dascrivere{la principessa che cresce lontano: com'è adesso}\par

\vspace{5mm}
\textit{A Zia Enza}\par
\vspace{2mm}
\dascrivere{quella che mette ordine nel caos}\par

\vspace{5mm}
\textit{A Zia Cristina}\par
\vspace{2mm}
\dascrivere{la paura delle catastrofi e il cuore immenso}\par

\vspace{5mm}
\textit{A mia nonna Maria}\par
\vspace{2mm}
A te, presenza costante in ogni cosa fin dal mio primo giorno, anzi, da prima ancora che nascessi. A te che, nonostante la lontananza, \dascrivere{continuazione di Claudio}\par

\vfill
\begin{flushright}
    \textit{Il vostro Claudio}
\end{flushright}

\clearpage
\vspace{10mm}
\begin{center}
    \textit{A Marta}
\end{center}
\vspace{5mm}

\dascrivere{nella triennale erano cinque anni, ora sei: cosa avete vissuto in questi due anni, cosa vuoi dirle adesso}\par

\clearpage
\vspace{10mm}
\begin{center}
    \textit{Ai miei amici}
\end{center}
\vspace{5mm}

\dascrivere{introduzione agli amici, «famiglia lontano da casa»: cosa è cambiato finita la triennale}\par

\vspace{5mm}
\textit{A Omar}\par
\vspace{2mm}
Ormai non viviamo più insieme e mi manca averti in casa. Ripenso a quando sei arrivato con quel valigione enorme rosso e io mi sono chiesto: «E mo' chi è questo?». Non immaginavo che, qualche anno dopo, mi avrebbe fatto così male trovare vuota la tua stanza.\par

È successo più di una volta: sono salito in automatico, senza pensarci, come avevo fatto tante altre volte. Ho aperto la porta e tu non c'eri. Per un attimo mi ero dimenticato che non vivevamo più insieme.\par

Con te in casa mi sentivo al sicuro. Sapevo che, se avevo bisogno di parlare, di sfogarmi o anche soltanto di stare un po' in compagnia, eri a una rampa di scale da me. Con te potevo dire quello che avevo in testa senza preparare un discorso, senza paura di essere giudicato. Potevo essere me stesso, anche nelle giornate in cui facevo fatica a sopportarmi.\par

Mi manca anche il rumore delle tue ciabatte mentre studiavamo, una di quelle cose a cui allora non facevo nemmeno caso e che adesso vorrei risentire. Mi manca entrare in camera tua per dirti una stupidaggine e finire a parlare per ore. Abbiamo condiviso le feste e le follie, ma oggi mi accorgo di quanto contassero anche quei momenti normali, quando bastava sapere che eri in casa. Mi manchi, Omar.\par

Nella triennale ti avevo chiamato il mio complice, mio fratello. Ormai siamo fratelli, con due caratteri completamente diversi che, in qualche modo, si completano. Anche se, quando ci mettiamo insieme, più che completarci combiniamo un caos. E mi manca pure quello.\par

Ora siamo lontani, ma continuo a sentirti vicino. Mi hai dimostrato più di una volta che ci tieni a me e che, se ho bisogno di qualcosa, mi basta scriverti: tu ci sei, mi ascolti e mi aiuti. Questo mi dà ancora quella sicurezza che sentivo quando abitavamo insieme.\par

Grazie per essermi rimasto vicino anche da lontano. Sei ancora la persona da cui so di poter andare, anche se adesso, per raggiungerti, non basta più una rampa di scale.\par

E comunque, adesso che lavori, aspetto ancora il mio viaggio all inclusive, tutto pagato. Dopo tutte queste belle parole, direi che me lo sono meritato.\par

\vspace{5mm}
\textit{A Cosimo}\par
\vspace{2mm}
\dascrivere{la prima cassa di birra, le feste del campus, e dopo}\par

\vspace{5mm}
\textit{A Gianmaria}\par
\vspace{2mm}
\dascrivere{chi è per te, un ricordo di questi anni}\par

\vspace{5mm}
\textit{A Youssef}\par
\vspace{2mm}
\dascrivere{rum, sigari, grigliate: cosa è venuto dopo}\par

\vspace{5mm}
\textit{A Rafa}\par
\vspace{2mm}
\dascrivere{l'estate sotto lo stesso tetto e quello che è seguito}\par

\vspace{5mm}
\textit{A Noemi}\par
\vspace{2mm}
\dascrivere{dal corso di inglese a oggi}\par

\vspace{5mm}
\textit{A Francesco}\par
\vspace{2mm}
Grazie ancora per ogni discussione, ogni risata, per la fiducia che mi dai e per la pazienza con cui reggi i nostri momenti di follia. Oggi le nostre strade ci portano in posti diversi, ma quello che c'è tra noi non si misura in chilometri: ogni volta che ci ritroviamo, tutto è esattamente dove l'abbiamo lasciato.\par

Non sempre è stato facile. Ci sono stati momenti in cui il nostro rapporto ha preso dei colpi, e avrebbe potuto rovinarsi. Se non è successo, è merito della volontà di entrambi, della fiducia che abbiamo l'uno nell'altro e della forza di dirci le cose in faccia, anche quando costava. È anche per questo che oggi so quanto vale quello che c'è tra noi.\par

Ti voglio bene, Ciccio, e sono davvero contento di averti conosciuto.\par

Ora che hai la patente nautica, però, ci aspettiamo un giro in barca. Io mi prenoto già, poi a fidarmi di te al timone ci penso quando siamo lì.\par

\vspace{5mm}
\textit{A Sofia}\par
\vspace{2mm}
La tua dolcezza e la tua sincerità sono un dono raro. Con le feste a sorpresa, la tua organizzazione, i pranzi e le cene che ci prepari, riesci sempre a farci sentire tutti a casa.\par

Hai un modo tutto tuo di prenderti cura delle persone: lo fai con le cose concrete, un piatto pronto, una sorpresa pensata nei dettagli, e lo fai senza mai farlo pesare. Grazie, Sofia, perché trasformi in casa ogni posto in cui siamo insieme.\par

Ti voglio un bene immenso. E la prossima volta che ci ritroviamo a tavola, tu cucini e io porto il dolce. Anche se, lo sappiamo, la più dolce resti sempre tu.\par

\vspace{5mm}
\dascrivere{chiusura agli amici}\par

\vfill
\begin{flushright}
    \textit{Claudio Dragotta}\\
    \textit{Roma, \the\year}
\end{flushright}

\clearpage
